% Sample tex provided by a group project for our Data Science course
% This is just a small project for a class & is unpublished

%% This is file `sample-authordraft.tex',
%% generated with the docstrip utility.
%%
%% The original source files were:
%%
%% samples.dtx  (with options: `authordraft')
%% 
%% IMPORTANT NOTICE:
%% 
%% For the copyright see the source file.
%% 
%% Any modified versions of this file must be renamed
%% with new filenames distinct from sample-authordraft.tex.
%% 
%% For distribution of the original source see the terms
%% for copying and modification in the file samples.dtx.
%% 
%% This generated file may be distributed as long as the
%% original source files, as listed above, are part of the
%% same distribution. (The sources need not necessarily be
%% in the same archive or directory.)
%%
%% Commands for TeXCount
%TC:macro \cite [option:text,text]
%TC:macro \citep [option:text,text]
%TC:macro \citet [option:text,text]
%TC:envir table 0 1
%TC:envir table* 0 1
%TC:envir tabular [ignore] word
%TC:envir displaymath 0 word
%TC:envir math 0 word
%TC:envir comment 0 0
%%
%%
%% The first command in your LaTeX source must be the \documentclass command.
\documentclass[sigconf,authordraft]{acmart}
%% NOTE that a single column version may required for 
%% submission and peer review. This can be done by changing
%% the \doucmentclass[...]{acmart} in this template to 
%% \documentclass[manuscript,screen]{acmart}
%% 
%% To ensure 100% compatibility, please check the white list of
%% approved LaTeX packages to be used with the Master Article Template at
%% https://www.acm.org/publications/taps/whitelist-of-latex-packages 
%% before creating your document. The white list page provides 
%% information on how to submit additional LaTeX packages for 
%% review and adoption.
%% Fonts used in the template cannot be substituted; margin 
%% adjustments are not allowed.

%%
%% \BibTeX command to typeset BibTeX logo in the docs
% \AtBeginDocument{%
%   \providecommand\
%   BibTeX{{%
%     \normalfont B\kern-0.5em{\scshape i\kern-0.25em b}\kern-0.8em\TeX}}}

%% Rights management information.  This information is sent to you
%% when you complete the rights form.  These commands have SAMPLE
%% values in them; it is your responsibility as an author to replace
%% the commands and values with those provided to you when you
%% complete the rights form.
\setcopyright{acmcopyright}
\copyrightyear{2023}
\acmYear{2018}
\acmDOI{XXXXXXX.XXXXXXX}

%% These commands are for a PROCEEDINGS abstract or paper.
\acmConference[Conference acronym 'XX]{Make sure to enter the correct
  conference title from your rights confirmation emai}{June 03--05,
  2018}{Woodstock, NY}
%
%  Uncomment \acmBooktitle if th title of the proceedings is different
%  from ``Proceedings of ...''!
%
%\acmBooktitle{Woodstock '18: ACM Symposium on Neural Gaze Detection,
%  June 03--05, 2018, Woodstock, NY} 
\acmPrice{15.00}
\acmISBN{978-1-4503-XXXX-X/18/06}


%%
%% Submission ID.
%% Use this when submitting an article to a sponsored event. You'll
%% receive a unique submission ID from the organizers
%% of the event, and this ID should be used as the parameter to this command.
%%\acmSubmissionID{123-A56-BU3}

%%
%% For managing citations, it is recommended to use bibliography
%% files in BibTeX format.
%%
%% You can then either use BibTeX with the ACM-Reference-Format style,
%% or BibLaTeX with the acmnumeric or acmauthoryear sytles, that include
%% support for advanced citation of software artefact from the
%% biblatex-software package, also separately available on CTAN.
%%
%% Look at the sample-*-biblatex.tex files for templates showcasing
%% the biblatex styles.
%%

%%
%% For managing citations, it is recommended to use bibliography
%% files in BibTeX format.
%%
%% You can then either use BibTeX with the ACM-Reference-Format style,
%% or BibLaTeX with the acmnumeric or acmauthoryear sytles, that include
%% support for advanced citation of software artefact from the
%% biblatex-software package, also separately available on CTAN.
%%
%% Look at the sample-*-biblatex.tex files for templates showcasing
%% the biblatex styles.
%%

%%
%% The majority of ACM publications use numbered citations and
%% references.  The command \citestyle{authoryear} switches to the
%% "author year" style.
%%
%% If you are preparing content for an event
%% sponsored by ACM SIGGRAPH, you must use the "author year" style of
%% citations and references.
%% Uncommenting
%% the next command will enable that style.
%%\citestyle{acmauthoryear}

%%
%% end of the preamble, start of the body of the document source.
\begin{document}

%%
%% The "title" command has an optional parameter,
%% allowing the author to define a "short title" to be used in page headers.
\title{Leveraging Convolutional Neural Networks for Accurate Material Detection in Mining Operations}

%%
%% The "author" command and its associated commands are used to define
%% the authors and their affiliations.
%% Of note is the shared affiliation of the first two authors, and the
%% "authornote" and "authornotemark" commands
%% used to denote shared contribution to the research.
\author{Pouria Rad}
\authornote{All authors contributed equally to this research.}
\email{peslamirad@augusta.edu}
% \orcid{1234-5678-9012}
% \author{G.K.M. Tobin}
% \authornotemark[1]
% \email{webmaster@marysville-ohio.com}
\affiliation{%
  \institution{Augusta University}
  \streetaddress{100 Grace Hopper Ln}
  \city{Augusta}
  \state{GA}
  \country{USA}
  \postcode{30901}
}

\author{Mahdieh Heidaripour}
\email{mheidaripour@augusta.edu}
\affiliation{%
  \institution{Augusta University}
  \streetaddress{100 Grace Hopper Ln}
  \city{Augusta}
  \state{GA}
  \country{USA}
  \postcode{30901}
}

\author{Asma Jodiri Akbarfam}
\email{ajodeiriakbarfam@augusta.edu}
\affiliation{%
  \institution{Augusta University}
  \streetaddress{100 Grace Hopper Ln}
  \city{Augusta}
  \state{GA}
  \country{USA}
  \postcode{30901}
}

\author{Seth Barrett}
\email{sebarrett@augusta.edu}
\affiliation{%
  \institution{Augusta University}
  \streetaddress{100 Grace Hopper Ln}
  \city{Augusta}
  \state{GA}
  \country{USA}
  \postcode{30901}
}

\author{Salil Verma}
\email{salverma@augusta.edu}
\affiliation{%
  \institution{Augusta University}
  \streetaddress{100 Grace Hopper Ln}
  \city{Augusta}
  \state{GA}
  \country{USA}
  \postcode{30901}
}

% \author{John Smith}
% \affiliation{%
%   \institution{The Th{\o}rv{\"a}ld Group}
%   \streetaddress{1 Th{\o}rv{\"a}ld Circle}
%   \city{Hekla}
%   \country{Iceland}}
% \email{jsmith@affiliation.org}

% \author{Julius P. Kumquat}
% \affiliation{%
%   \institution{The Kumquat Consortium}
%   \city{New York}
%   \country{USA}}
% \email{jpkumquat@consortium.net}

%%
%% By default, the full list of authors will be used in the page
%% headers. Often, this list is too long, and will overlap
%% other information printed in the page headers. This command allows
%% the author to define a more concise list
%% of authors' names for this purpose.
\renewcommand{\shortauthors}{Trovato and Tobin, et al.}

%%
%% The abstract is a short summary of the work to be presented in the
%% article.
\begin{abstract}
This paper presents a comprehensive approach to rock quality designation (RQD) estimation using object detection techniques. The proposed method leverages Faster R-CNN architecture to accurately identify and label rock and wood pieces in images. We introduce a dataset of annotated box images, each containing various types of rocks and woods. Through the implementation of two separate Faster R-CNN models, we effectively detect and analyze these objects, enabling the calculation of RQD values for different runs within the dataset. Additionally, we detail the training process, including hyperparameter optimization and data augmentation strategies, for both the rock and wood models. Our experimental results demonstrate the effectiveness of the proposed method in accurately predicting RQD groups, showcasing the potential for enhanced geotechnical analysis and mineral exploration.
\end{abstract}

%%
%% The code below is generated by the tool at http://dl.acm.org/ccs.cfm.
%% Please copy and paste the code instead of the example below.
%%
\begin{CCSXML}
<ccs2012>
 <concept>
  <concept_id>00000000.0000000.0000000</concept_id>
  <concept_desc>Do Not Use This Code, Generate the Correct Terms for Your Paper</concept_desc>
  <concept_significance>500</concept_significance>
 </concept>
 <concept>
  <concept_id>00000000.00000000.00000000</concept_id>
  <concept_desc>Do Not Use This Code, Generate the Correct Terms for Your Paper</concept_desc>
  <concept_significance>300</concept_significance>
 </concept>
 <concept>
  <concept_id>00000000.00000000.00000000</concept_id>
  <concept_desc>Do Not Use This Code, Generate the Correct Terms for Your Paper</concept_desc>
  <concept_significance>100</concept_significance>
 </concept>
 <concept>
  <concept_id>00000000.00000000.00000000</concept_id>
  <concept_desc>Do Not Use This Code, Generate the Correct Terms for Your Paper</concept_desc>
  <concept_significance>100</concept_significance>
 </concept>
</ccs2012>
\end{CCSXML}

\ccsdesc[500]{Do Not Use This Code~Generate the Correct Terms for Your Paper}
\ccsdesc[300]{Do Not Use This Code~Generate the Correct Terms for Your Paper}
\ccsdesc{Do Not Use This Code~Generate the Correct Terms for Your Paper}
\ccsdesc[100]{Do Not Use This Code~Generate the Correct Terms for Your Paper}

%%
%% Keywords. The author(s) should pick words that accurately describe
%% the work being presented. Separate the keywords with commas.
\keywords{Object Detection,
Rock Quality Designation (RQD), R-CNN}

%% A "teaser" image appears between the author and affiliation
%% information and the body of the document, and typically spans the
%% page.
% \begin{teaserfigure}
%   \includegraphics[width=\textwidth]{sampleteaser}
%   \caption{Seattle Mariners at Spring Training, 2010.}
%   \Description{Enjoying the baseball game from the third-base
%   seats. Ichiro Suzuki preparing to bat.}
%   \label{fig:teaser}
% \end{teaserfigure}

\received{20 February 2007}
\received[revised]{12 March 2009}
\received[accepted]{5 June 2009}

%%
%% This command processes the author and affiliation and title
%% information and builds the first part of the formatted document.
\maketitle

\section{Introduction}

Geotechnical engineers play a critical role in ensuring the stability and safety of infrastructure projects involving rock formations. Understanding the mechanical properties and integrity of rock masses is paramount in designing foundations, tunnels, slopes, and other structures. Rock Quality Designation (RQD) serves as a key parameter in this assessment, providing insights into the proportion of intact rock versus fractured material in a core sample. For geotechnical engineers, a high RQD value indicates a greater presence of sound rock, suggesting enhanced structural stability and suitability for construction. Conversely, a low RQD value signifies a higher occurrence of fractures, indicating potential challenges in stability and strength. Traditionally, RQD assessment involves manual inspection and measurement of core samples, a process susceptible to subjectivity and resource-intensive efforts. A shift towards leveraging Machine Learning (ML) techniques is imperative, as it introduces efficiency, objectivity, and scalability to the evaluation process, ensuring more accurate and timely assessments for geotechnical projects.

The conventional approach to RQD assessment, relying on manual inspection, poses inherent challenges. Human subjectivity, variability in interpretation, and the labor-intensive nature of the process can lead to inconsistencies in results. Additionally, the traditional method may be insufficient in handling large datasets, limiting its applicability in projects with extensive drilling activities. The integration of ML technologies offers a transformative solution to these challenges. By automating the identification of intact rock and fractures in core samples, ML algorithms enhance the accuracy and speed of RQD calculations. This not only streamlines the workflow for geotechnical engineers but also introduces a more objective and standardized approach. The implementation of ML in RQD estimation aligns with the industry's demand for advanced technologies that can improve decision-making, reduce manual labor, and enhance the overall reliability of geological assessments for infrastructure development and mining activities.

In this paper, we present an innovative approach to identify and measure the Rock Quality Designation (RQD) of geological core samples. The RQD is a crucial parameter for characterizing rock quality and is vital in various geological and mining applications \cite{1,2,7}. Our method leverages object detection to automatically label and measure rocks and wood pieces within core box images, drawing upon recent advances in the application of convolutional neural networks (CNNs) and object detection techniques in similar contexts \cite{1,3,20}.




\section{Background}
Convolutional Neural Networks (CNNs) have significantly advanced the field of image recognition and analysis, particularly in specialized areas like geological assessment and material identification. The Rock Quality Designation (RQD) is a important measure in geology, especially for assessing the quality of rock. Recent progress in CNNs has enabled more precise and efficient estimation of RQD. Studies such as those by \cite{4}, \cite{5}, and \cite{6} have established the foundation for using CNNs in intricate image analysis tasks, which is directly pertinent to geological core sample analysis and RQD analysis. \\

The work by \cite{7} and \cite{8} further demonstrates the real-wrold application of CNNs in geological contexts, underscoring the potential of these models to transform conventional methods of rock quality assessment and RQD calculation. The use of CNNs for image classification, as discussed in \cite{9}, provides vital understandings into fine-tuning these networks for specific assignments like RQD estimation in rock. \\

Furthermore, Faster R-CNN, an evolution of CNNs, has shown significant possibility in material labeling, improving speed and accuracy over traditional methods. Papers such as \cite{10}, \cite{11}, and \cite{12} discuss advances in object detection using Faster R-CNN, demonstrating its capability in accurately classifying and identifying diverse materials. The works of \cite{13} and \cite{14} specifically highlight the application of Faster R-CNN in complex environments, which can be comparable to material identification in mining operations. \\

Enhancements in Faster R-CNN models, as detailed in \cite{15} and \cite{16}, provide a framework for improved detection methods in industrial and environmental contexts. As explored in \cite{17}, \cite{18}, and \cite{1}, the advancements indicate the potential for applying Faster R-CNN in the precise classification and labeling of geological materials. \\

\section{Related Work}
Jiang et al.\cite{31} developed an application using Faster R-CNN deep learning for automated identification and counting of geological features in drill core images. They created a dataset of core sample images with labels for joints, fracture zones, round trip marks, etc. A Faster R-CNN model was trained to identify and localize the geological features in the images. In testing on real core sample folders, the AI model achieved over 90\% accuracy in feature detection. The overall application with batch processing was 48 times faster than manual identification. This demonstrates the potential of deep learning techniques like Faster R-CNN to accurately automate the detection and quantification of materials, textures, and objects in images from mining contexts. The specialized mining image dataset and speed improvement validate the utility of such approaches for material characterization in mining. \\

Alzubaidi et al.\cite{1} proposes a convolutional neural network (CNN) method for automated rock characterization and calculation of the rock quality designation (RQD) from drill core images. The CNN model detects intact core fragments and excludes fractures, crushed cores, empty tray areas, and non-rock objects to estimate RQD. The model was trained on 6400 sandstone core images and tested on sandstone and limestone cores, achieving average error rates of 2.58\% and 3.17\% respectively compared to manual RQD calculations. The CNN-based approach provides fast, robust RQD estimation directly from images, eliminating subjective manual measurements and enabling automated analysis of drill core databases. \\

Fernández-Alonso et al.\cite{32} developed a convolutional neural network (CNN) system to automatically identify Roman mining remains from aerial drone images. They created a dataset of drone images showing mining features and implemented a CNN architecture to extract features and classify the images. In testing, the CNN achieved approximately 95\% accuracy in identifying channels, roads, and mines, outperforming other techniques. The authors demonstrate the potential of deep learning and drone imagery to automate detection of materials and objects in mining. The specialized dataset and high accuracy highlight CNNs as a promising approach for mining feature detection without human subjectivity. \\

% Same as 2 paragraphs above
% Alzubaidi et al.\cite{1} proposes a convolutional neural network (CNN) method for automated rock characterization and calculation of the rock quality designation (RQD) from drill core images. The CNN model detects intact core fragments and excludes fractures, crushed cores, empty tray areas, and non-rock objects to estimate RQD. The model was trained on 6400 sandstone core images and tested on sandstone and limestone cores, achieving average error rates of 2.58\% and 3.17\% respectively compared to manual RQD calculations. The CNN-based approach provides fast, robust RQD estimation directly from images, eliminating subjective manual measurements and enabling automated analysis of drill core databases. \\

\section{Motivation}
%Accurate and efficient RQD assessment is essential for geological analysis. The conventional manual approach is time-consuming and prone to human errors \cite{2,7,19,22}. Our motivation is to develop an automated system that can reliably determine the RQD value from images, significantly improving productivity in geological assessments \cite{1,3,20}. \\
Accurate and efficient Rock Quality Designation (RQD) assessment stands as a cornerstone in geological analysis, crucial for understanding the integrity of rock masses. The conventional manual approach to RQD assessment is not only time-consuming but is also susceptible to human errors \cite{2, 7, 19, 22}. Geotechnical engineers and geological researchers encounter challenges in consistently and objectively measuring the proportion of intact rock versus fractured material within core samples using traditional methods. These challenges underline the critical need for an innovative and automated solution.

Our motivation stems from the inherent limitations of the conventional manual approach and the imperative to modernize and streamline RQD assessment processes \cite{1, 3, 20}. The intricacies involved in manually inspecting and measuring core samples often lead to variations in results, introducing subjectivity and potentially compromising the accuracy of geological assessments. To address these challenges, we aspire to develop an automated system capable of reliably determining RQD values from images \cite{1, 3, 20}. The primary goal is to significantly enhance productivity in geological assessments, offering a more accurate, efficient, and objective means of evaluating rock quality. This motivation aligns with the broader industry trend of adopting advanced technologies, such as Machine Learning (ML), to revolutionize traditional geological methodologies \cite{1, 3, 20}. Through this endeavor, we aim to contribute to the ongoing advancements in geotechnical analysis, providing geologists and engineers with a more robust and automated tool for RQD estimation.

\section{Proposed Method}
\subsection{Data Preparation}

Our dataset comprises core box images containing rock and wood pieces. Each object in the images is labeled using white rectangles for identification. This annotation approach is consistent with standard practices in object detection tasks using CNNs \cite{10,13,15}.

\begin{figure}[h]
  \centering
  \includegraphics[width=0.4\textwidth]{image/box-annotations.jpg}
  \caption{Sample Object Annotations}
\end{figure}
\begin{figure}[h]
  \centering
  \includegraphics[width=0.4\textwidth]{image/object-dims.jpg}
  \caption{Object Parameters}
\end{figure}
We provide an example of the parameters associated with a labeled object. Such detailed annotation is crucial for the accurate training of machine learning models, as evidenced in similar studies \cite{11,14,16}.

The dataset includes essential information stored in an Excel file:
\begin{center}
\begin{tabular}{|c|p{0.3\textwidth}|}
  \hline
  \textbf{Column name} & \textbf{Description} \\
  \hline
  \texttt{image\_name} & ID of a photo (box) \\
  \texttt{label\_name} & The name of the detected object \\
  \texttt{xmin} & x coordinate of the upper-left corner \\
  \texttt{ymin} & y coordinate of the upper-left corner \\
  \texttt{width} & Length of object rectangle (in x axis) \\
  \texttt{height} & Height of object rectangle (on y axis) \\
  \texttt{image\_width} & Image width in pixels \\
  \texttt{image\_height} & Image height in pixels \\
  \hline
\end{tabular}
\end{center}

\subsection{RQD Calculation}

RQD, a key metric for assessing rock quality, is calculated as the percentage of fractures in the rock mass. The RQD value for each run is determined based on the presence of rocks larger than 10 cm. Each run is defined between vertical wood pieces in the core box, following the standard approach outlined in existing studies \cite{1,2,7,19,22}. \\

For example, in a run, the RQD is calculated as follows:

\begin{figure}[h]
  \centering
  \includegraphics[width=0.3\textwidth]{image/rqd-example.png}
  \caption{RQD Calculation Example}
\end{figure}
We categorize RQD values into groups:

\begin{table}[H]
\centering
\begin{tabular}{|c|c|}
  \hline
  \textbf{RQD percent range} & \textbf{Class} \\
  \hline
  0 -- 25 & 1 \\
  25 -- 50 & 2 \\
  50 -- 75 & 3 \\
  75 -- 90 & 4 \\
  90 -- 100 & 5 \\
  \hline
\end{tabular}
\end{table}

The length of the core box is standardized at 1.1 meters along the x-axis, with purple-colored wood markers, a method that aligns with practices described in \cite{3,20,21} for consistent and accurate measurement in geological assessments.


\section{Model Configuration}

To address the object detection challenge, we employed the Faster R-CNN architecture \cite{13,14,26}. We developed two separate models, one for identifying rocks and another for identifying wood pieces. Our modeling process included optimizing key hyperparameters: 

\textit{Learning rate, Momentum, Weight decay, Batch size.}

Hyperparameter optimization was performed using hyperOpt \cite{16,17,18}, resulting in the following hyperparameters for our models:

\textbf{Rock Model Hyperparameters}:
  - Batch size: 8
  - Learning rate: 0.0084
  - Momentum: 0.9445
  - Weight decay: 0.0074

\textbf{Wood Model Hyperparameters}:
  - Batch size: 8
  - Learning rate: 0.0060
  - Momentum: 0.9708
  - Weight decay: 0.0015

We implemented early stopping with patience set to 4 and a minimum delta of 0.01. This approach aligns with best practices in machine learning to prevent overfitting \cite{25,27,29}.

The learning curves of both models are shown in Figures \ref{fig:curveR} and \ref{fig:curveW}
\begin{figure}[h]
  \centering
  \includegraphics[width=0.4\textwidth]{image/training_curve_rock_model.png}
  \caption{Training Curve Rock Model.}
  \label{fig:curveR}
\end{figure}

\begin{figure}[h]
  \centering
  \includegraphics[width=0.4\textwidth]{image/training_curve_wood_model.png}
  \caption{Training Curve Wood Model.}
  \label{fig:curveW}
\end{figure}


\section{Evaluation and Results}
We evaluated the proposed method on a comprehensive dataset of diverse box images, utilizing the Faster R-CNN architecture for object detection and RQD estimation. The evaluation was conducted on Google Colab, and the project was successfully executed using Python 3.9.16. The project dependencies were carefully specified, with the following package versions utilized:
numpy==1.22.4, tqdm==4.65.0, torchvision==0.14.1+cu116, torch==1.13.1+cu116, scipy==1.10.1, hyperopt==0.2.7, and matplotlib==3.7.1. \\

\subsection{Dataset}

Our training dataset is stored in the \texttt{train} folder, while the test dataset is in the \texttt{test-rqd} folder. The depths of runs for the test dataset are obtained from \texttt{from-to-rqd.xlsx}, as per established procedures in similar studies \cite{5,20,24}. \\

The structure of a RunId is based on "mineName-boreholeName-photoNumber-runNumber," with photo and run numbers starting from 1, aligning with data organization practices in machine learning \cite{11,14,16}. \\

\subsection{Metric}

We evaluate our model's performance using the accuracy criterion, which is multiplied by 1000 to calculate the score. The test data evaluation is initially based on 30\% of the data to prevent overfitting, following the guidelines in \cite{25,27,29}. After the competition, the entire test dataset is used for the final update of the scoreboard. \\

\subsection{Results}

We require your model predictions on the test dataset to be saved in an \texttt{output.csv} file. This file should include two columns: \texttt{RunId} and \texttt{Prediction}. Each row should contain the corresponding \texttt{RunId} and your predicted RQD group.

We used 15 percent of the dataset for validation and the remaining 85 percent for training, consisting of 156 samples. Augmentations and transformations, such as horizontal flips, were applied to improve model robustness, a technique supported by \cite{10,13,15}. \\

Learning was conducted over 40 epochs with early stopping, an approach that aligns with best practices in machine learning to prevent overfitting \cite{25,27,29}. \\

\subsection{Prediction and Validation}

We applied the trained models to a sample from the validation dataset and observed the model's outputs. \\

\textbf{Rock Model Output}:

\begin{figure}[h]
  \centering
  \includegraphics[width=0.4\textwidth]{image/rock_model_sample_org.png}
  \caption{Rock Model Output - Original}
\end{figure}

\begin{figure}[h]
  \centering
  \includegraphics[width=0.4\textwidth]{image/rock_model_sample_pred.png}
  \caption{Rock Model Output - Predicted}
\end{figure}

\textbf{Wood Model Output}:

\begin{figure}[h]
  \centering
  \includegraphics[width=0.4\textwidth]{image/wood_model_sample_org.png}
  \caption{Wood Model Output - Original}
\end{figure}

\begin{figure}[h]
  \centering
  \includegraphics[width=0.4\textwidth]{image/wood_model_sample_pred.png}
  \caption{Wood Model Output - Predicted}
\end{figure}


\section{Conclusion}

In conclusion, our approach leverages object detection to automate RQD assessment in geological core samples. By employing Faster R-CNN models and rigorous hyperparameter optimization \cite{13,14,16,26}, we've demonstrated promising results. This method not only enhances efficiency but also ensures accuracy in geological evaluations, echoing findings from studies like \cite{1,3,20}. We encourage further exploration and application of this technique in geology and mining scenarios, as it holds great potential for revolutionizing traditional methods \cite{2,7,19,22}. \\


\bibliographystyle{acm}
\bibliography{main}


\end{document}
\endinput
%%
%% End of file `sample-authordraft.tex'.
